\documentclass{article}
\usepackage[utf8]{inputenc}
\usepackage{a4wide}

\begin{document}

\section*{ซวยวันสอบ}

ในวันสอบ สอวน ค่าย 2 วิชาคอมพิวเตอร์ ผู้ปกครองของดานาไม่สะดวกมาส่งดานาที่มหาวิทยาลัย

ทำให้ดานาต้องเดินทางมาสอบที่มหาวิทยาลัยด้วยตัวเอง

แต่บ้านของดานาก็อยู่ห่างไกลจากมหาวิทยาลัยมากเหลือเกิน จริง ๆ ก็อยากนั่ง Grab แต่ net

ดันหมดพอดี หรือ นั่ง Taxi แต่เป็นช่วงต้นเดือนเงินสดไม่พอเพราะแม่ลืมให้เงินทิ้งไว้

ดังนั้นเธอจึงต้องขึ้นรถประจำทางเพื่อมาที่มหาวิทยาลัยเท่านั้น เส้นทางจากบ้านไปมหาวิทยาลัยนั้น

ดานาลองนึกวาดภาพในจินตนาการก็เหมือนกับแทนด้วยส่วนของเส้นตรง

โดยบนส่วนของเส้นตรงนี้มีจุดจอดรถประจำทางทั้งหมด n + 1 จุด

จุดจอดรถประจำทางแต่ละจุดมีหมายเลขตั้งแต่ 0 ถึง n ซึ่งเรียงตามลำดับที่จอดจากบ้านของดานา

โดยจุดจอดที่บ้านของดานามีหมายเลข 0 และจุดจอดที่มหาวิทยาลัยมีหมายเลข n



มีรถประจำทางทั้งหมด m คันที่วิ่งระหว่างบ้านกับมหาวิทยาลัย รถประจำทางคันที่ i

จะไปจากจุดจอด s{[}i{]} ไปยัง t{[}i{]} (s{[}i{]} \textless{} t{[}i{]})

โดยจะผ่านจุดจอดระหว่างทางที่เรียงตามลำดับตามเส้นตรงนี้

ดานาเองไม่คุ้นเคยกับการขึ้นรถประจำทาง

ดังนั้น\ul{\textbf{เมื่อเธอเลือกขึ้นรถคันไหนเธอจะไม่ลงจากรถรถประจำทางคนนั้นจนกว่าจะสามารถนั่งรถไปได้ใกล้มหาวิทยาลัยมากที่สุด}}นั่นคือ

ดานาสามารถขึ้นรถประจำทางคันที่ i ที่จุดจอดใดก็ได้ระหว่าง s{[}i{]} ถึง t{[}i{]} - 1

แต่เธอจะลงจากรถประจำทางคันนั้นได้ที่จุดจอด t{[}i{]} เท่านั้น



เนื่องจากเป็นเช้าวันสอบดานาเองต้องการประหยัดพลังงานมากที่สุด\ul{\textbf{เธอจะไม่เดินระหว่างป้ายรถประจำทาง

และเธอก็ไม่เดินย้อนกลับในทิศทางจากมหาวิทยาลัยมาบ้านด้วยเช่นกัน}}



ดานาต้องการรู้ว่ามีกี่วิธีที่เธอจะไปถึงมหาวิทยาลัยได้

โดยสองวิธีจะถือว่าแตกต่างกันถ้าดานาเดินทางข้ามบางช่วงระหว่างจุดจอดรถประจำทางบนรถประจำทางที่แตกต่างกัน

จำนวนวิธีอาจจะเยอะมาก\textbf{\ul{ดังนั้นคำตอบจะเป็นเศษของการหาค่าผลลัพธ์ที่ได้จากการหารด้วย

1000000007}}



\textbf{Input:}บรรทัดแรกประกอบด้วยจำนวนเต็มสองจำนวน n และ m (1 ≤ n ≤

10\^{}9, 0 ≤ m ≤ 10\^{}5)



จากนั้นมี m บรรทัดที่ประกอบด้วยจำนวนเต็มสองจำนวน s{[}i{]}, t{[}i{]}

ซึ่งเป็นหมายเลขของจุดจอดเริ่มต้นและจุดจอดสิ้นสุดของรถประจำทาง (0 ≤ s{[}i{]}

\textless{} t{[}i{]} ≤ n)



\textbf{Output:}เศษจากการเอาจำนวนวิธีการเดินทางไปถึงมหาวิทยาลัยหารด้วย

1000000007



\paragraph{\texorpdfstring{\textbf{ตัวอย่าง}}{ตัวอย่าง}}\label{uxe15uxe27uxe2duxe22uxe32uxe07}



{\def\LTcaptype{none} % do not increment counter

\begin{longtable}[]{@{}

  >{\raggedright\arraybackslash}p{(\linewidth - 4\tabcolsep) * \real{0.3333}}

  >{\raggedright\arraybackslash}p{(\linewidth - 4\tabcolsep) * \real{0.3333}}

  >{\raggedright\arraybackslash}p{(\linewidth - 4\tabcolsep) * \real{0.3333}}@{}}

\toprule\noalign{}

\begin{minipage}[b]{\linewidth}\centering

Input\\

\strut

\end{minipage} & \begin{minipage}[b]{\linewidth}\centering

Output\\

\strut

\end{minipage} & \begin{minipage}[b]{\linewidth}\centering

คำอธิบาย\\

\strut

\end{minipage} \\

\midrule\noalign{}

\endhead

\bottomrule\noalign{}

\endlastfoot

\begin{minipage}[t]{\linewidth}\raggedright

3 2\\

0 2\\

2 3\\

\strut

\end{minipage} & \begin{minipage}[t]{\linewidth}\raggedright

1\\

\strut

\end{minipage} & \begin{minipage}[t]{\linewidth}\raggedright

ในกรณีนี้ มีเพียงทางเดียวที่ดานาจะไปถึงโรงเรียน: ขึ้นรถประจำทางหมายเลข 1

ที่จุดจอดหมายเลข 0 ไปลงที่จุดจอดหมายเลข 2 จากนั้นขึ้นรถประจำทางหมายเลข 2

จากจุดจอดหมายเลข 2 ไปลง\\

ที่จุดจอดหมายเลข 3 ซึ่งคือมหาวิทยาลัย\\

\strut

\end{minipage} \\

\begin{minipage}[t]{\linewidth}\raggedright

3 3\\

0 1\\

1 2\\

0 2\\

\strut

\end{minipage} & \begin{minipage}[t]{\linewidth}\raggedright

0\\

\strut

\end{minipage} & \begin{minipage}[t]{\linewidth}\raggedright

ในกรณีนี้ไม่มีรถบัสคันไหนที่ไปถึงจุดจอดหมายเลข 3 ซึ่งเป็นจุดจอดที่มหาวิทยาลัย ดังนั้นคำตอบคือ

0\\

\strut

\end{minipage} \\

\begin{minipage}[t]{\linewidth}\raggedright

7 7\\

0 1\\

0 2\\

0 3\\

0 4\\

0 5\\

0 6\\

0 7\\

\strut

\end{minipage} & \begin{minipage}[t]{\linewidth}\raggedright

64\\

\strut

\end{minipage} & \begin{minipage}[t]{\linewidth}\raggedright

ในกรณีนี้ ดานาสามารถเลือกขึ้นรถประจำทางคันไหนก็ได้จากจุดจอดหมายเลข 0

ไปยังจุดจอดหมายเลข 7 ซึ่งเป็นที่มหาวิทยาลัยได้

โดยสามารถขึ้นรถประจำทางจากจุดจอดใดก็ได้ในช่วงก่อนถึงจุดจอดหมายเลข 7 ดังนั้นคำตอบคือ

2\^{}6 = 64\\

\strut

\end{minipage} \\

\end{longtable}

}



\subsection*{Input}
บรรทัดแรก:

จำนวนเต็มสองจำนวน\texttt{n}และ\texttt{m}โดย\texttt{n}คือหมายเลขของจุดจอดสุดท้าย

(ซึ่งเป็นที่ตั้งของมหาวิทยาลัย)\texttt{m}คือจำนวนรถประจำทางทั้งหมด(1 ≤ n ≤ 10\^{}9,

0 ≤ m ≤ 10\^{}5)



จากนั้นมี\texttt{m}บรรทัด

โดยแต่ละบรรทัดประกอบด้วยจำนวนเต็มสองจำนวน\texttt{s{[}i{]}}และ\texttt{t{[}i{]}}โดย\texttt{s{[}i{]}}คือหมายเลขจุดจอดต้นทางของรถประจำทางคันที่\texttt{i}และ\texttt{t{[}i{]}}คือหมายเลขจุดจอดปลายทางของรถประจำทางคันที่\texttt{i}(0

≤ s{[}i{]} \textless{} t{[}i{]} ≤ n)



\subsection*{Output}
{เศษจากการเอาจำนวนวิธีการเดินทางไปถึงมหาวิทยาลัยหารด้วย 1000000007}\\



\end{document}
